% `tex.web` §1094: `end_graf` is about the mode, and only one mode.
%
%     procedure end_graf;
%     begin if mode=hmode then
%       begin if head=tail then pop_nest {null paragraphs are discarded}
%       else line_break(widow_penalty);
%       normal_paragraph;
%       error_count:=0;
%       end;
%     end;
%
% `mode=hmode`, not `abs(mode)=hmode`. A restricted horizontal list — the
% inside of an \hbox — is not a paragraph, so there is nothing to break and
% nothing for §1070's `normal_paragraph` to reset.
%
% We called it whatever the mode was. That is invisible until something has
% been set that `normal_paragraph` clears: inside `{\looseness 5 …}` an \hbox
% reset \looseness to 0 within the *box's own* group, and the box's closing
% brace gave it straight back — a `{restoring \looseness=5}` TeX does not
% write, and, worse, a \looseness of 0 in force for the whole of the box.
%
% Which is also why the fix has to leave `\par` alone. §1094 puts `\par` in the
% big switch twice, and only the horizontal one is `end_graf`:
%
%     vmode+par_end: begin normal_paragraph;
%       if mode>0 then build_page; end;
%     hmode+par_end: begin if align_state<0 then off_save; end_graf;
%       if mode=vmode then build_page; end;
%
% so a `\par` where there is no paragraph still resets the shape, and a `\par`
% in restricted horizontal mode does nothing whatever.
\catcode`\{=1 \catcode`\}=2
\tracingonline=1 \tracingcommands=2 \tracingrestores=1
\hsize=100pt \parindent=0pt
\font\rm=cmr10 \rm

% The \hbox is inside the group, so anything it resets is restored at *its*
% brace and shows up as a line of its own.
{\looseness 5 \hangindent 3pt A\hbox{}B\par}

% A \vbox does call `normal_paragraph` — as it opens (§1083), not as it
% closes — so this one restore is TeX's too.
{\looseness 5 \vbox{}}

% `\par` with no paragraph under way: the shape is reset, and the restore that
% follows is the group's.
{\looseness 5 \hangafter 4 \par}

% ...and `\par` in restricted horizontal mode, which is neither.
{\looseness 5 \setbox0\hbox{A\par B}}
\showthe\looseness
\end
